\section{Related work}
\label{fexpr-comparison}

\begin{note}
% claude:PR-STYLE-0002
% MTODO: use "fexprs", "fexpr-objects" everywhere
The concept still appears as ``F-expressions'', ``FExpr''/``FExprs''/
``FExpr-objects'' and ``fexprs''. Use one form throughout, e.g. ``fexpr'' and
``fexprs'' as in the cited literature, including the table and the footnote.
\end{note}

The techniques applied in this paper achieve results similar to systems using
F-expressions~\cite{fexpr-shutt,fexpr-pe}.

F-expressions are a different kind of metaprogramming mechanism that
unifies functions and macros as \emph{FExpr}-objects%
\footnote{
    The exact definition of the term \emph{FExpr} varies throughout
    literature. Here it is used in the sense of applicative object
    that may choose whether to have its operands evaluated before
    application~\cite{fexpr-shutt}.
}, which, depending on the implementation,
can be first-class and as flexible as unhygienic macros,%
~\cite{kohlbecker1986syntactic,fexpr-shutt} and can be implemented
in a performant way~\cite{fexpr-pe}. The core mechanism of FExprs is
\paraphrase{
to tag the object with evaluation metadata and consult that metadata
when evaluating a combination to decide whether to evaluate the operands}.
FExpr-objects are first-class regardless of their evaluation semantics.

Similar to modern FExpr-based systems, ILD's macro expansion facility matches
unhygienic LISP macros in expressive power and \paraphrase{provides first-class
status and runtime semantics with no separate expansion phase for macros}.
\begin{note}
% claude:PR-TECHNICAL-0001, codex:PR-TECHNICAL-0011
% MTODO: add a line for Kernel in the table
The claim that no language with first-class fexprs supports first-class
continuations is false: Kernel has first-class operatives and \ild{call/cc}
(Shutt, Revised\(^{-1}\) Report on the Kernel Programming Language,
Section 7). Remove or narrow the claim and compare ILD's callsite-triggered
expansion with Kernel's operative/applicative mechanism instead.
\end{note}
We argue that callsite-triggered macro expansion is simpler to implement and
reason about while being as expressive.
